% GENERATED FILE. Edit canonical lessons, then run npm run generate:books.
% canonical-source-hash: fnv1a64:fbae8eff6926ff46
% canonical-lessons: MR-C131-padda, MR-C131-kulup, MR-C131-badli, MR-C131-dori, MR-C131-khila

\chapter{Curtain, Lock, Bucket, Rope, Nail}
\label{ch:mr-curtain-lock-bucket-rope-nail}
\hlchaptermodality{131}

\begin{chapteropening}
\textbf{By the end of this chapter:} \emph{I can say a curtain, a lock, a bucket, a rope and a nail.}

Complete the last lesson of chapter 131: I can say a curtain, a lock, a bucket, a rope and a nail.
\end{chapteropening}

\section[\texorpdfstring{\mr{पडदा} (paḍdā)}{पडदा (paḍdā)}]{\mr{पडदा} (paḍdā) — a curtain}

\label{lesson:MR-C131-padda}



\begin{quote}
Before the new one: say the Marathi for a bedsheet, then the Marathi for a mattress.
\end{quote}

\subsection*{You'll want to know: \mr{पडदा}}
\textbf{\mr{पडदा}} — \emph{paḍdā} — \textquotedblleft{}a curtain\textquotedblright{}.

\begin{cousinweb}
\textbf{\mr{पडदा}} (\emph{paḍdā}) is a loan, usually traced to Persian \emph{parda}, \textquotedblleft{}curtain, veil\textquotedblright{}.
\end{cousinweb}

\begin{grammarlens}[title={What you've built}]
One more word for this chapter.
\end{grammarlens}

\subsection*{Your turn}
\begin{itemize}
\raggedright
  \item \emph{Say it:} \emph{paḍdā}
  \item \emph{Say it:} \emph{paḍdā}, once more
  \item \emph{Say it:} say \emph{paḍdā}
  \item \emph{Recall:} say the Marathi for a bedsheet, then the Marathi for a mattress, then say \emph{paḍdā} again
\end{itemize}

\begin{tcolorbox}[breakable,colback=teal!4,colframe=teal!35!black,title={Before you move on}]
What does \mr{पडदा} mean? (\textquotedblleft{}A curtain\textquotedblright{}.) Say it once more.
\end{tcolorbox}
\section[\texorpdfstring{\mr{कुलूप} (kulūp)}{कुलूप (kulūp)}]{\mr{कुलूप} (kulūp) — a lock}

\label{lesson:MR-C131-kulup}



\begin{quote}
Before the new one: say the Marathi for a mattress, then the Marathi for a curtain.
\end{quote}

\subsection*{You'll want to know: \mr{कुलूप}}
\textbf{\mr{कुलूप}} — \emph{kulūp} — \textquotedblleft{}a lock\textquotedblright{}.

\begin{cousinweb}
\textbf{\mr{कुलूप}} (\emph{kulūp}) is a loan, usually traced through Persian \emph{qulf} to Arabic \emph{qufl}, \textquotedblleft{}lock\textquotedblright{}, whose root is \emph{q-f-l} (Lane).
\end{cousinweb}

\begin{grammarlens}[title={What you've built}]
One more word for this chapter.
\end{grammarlens}

\subsection*{Your turn}
\begin{itemize}
\raggedright
  \item \emph{Say it:} \emph{kulūp}
  \item \emph{Say it:} \emph{kulūp}, once more
  \item \emph{Say it:} say \emph{kulūp}
  \item \emph{Recall:} say the Marathi for a mattress, then the Marathi for a curtain, then say \emph{kulūp} again
\end{itemize}

\begin{tcolorbox}[breakable,colback=teal!4,colframe=teal!35!black,title={Before you move on}]
What does \mr{कुलूप} mean? (\textquotedblleft{}A lock\textquotedblright{}.) Say it once more.
\end{tcolorbox}
\section[\texorpdfstring{\mr{बादली} (bādlī)}{बादली (bādlī)}]{\mr{बादली} (bādlī) — a bucket}

\label{lesson:MR-C131-badli}



\begin{quote}
Before the new one: say the Marathi for a curtain, then the Marathi for a lock.
\end{quote}

\subsection*{You'll want to know: \mr{बादली}}
\textbf{\mr{बादली}} — \emph{bādlī} — \textquotedblleft{}a bucket\textquotedblright{}.

\begin{cousinweb}
\textbf{\mr{बादली}} (\emph{bādlī}) is a loan, usually traced to Portuguese \emph{balde}, \textquotedblleft{}bucket\textquotedblright{}.
\end{cousinweb}

\begin{grammarlens}[title={What you've built}]
One more word for this chapter.
\end{grammarlens}

\subsection*{Your turn}
\begin{itemize}
\raggedright
  \item \emph{Say it:} \emph{bādlī}
  \item \emph{Say it:} \emph{bādlī}, once more
  \item \emph{Say it:} say \emph{bādlī}
  \item \emph{Recall:} say the Marathi for a curtain, then the Marathi for a lock, then say \emph{bādlī} again
\end{itemize}

\begin{tcolorbox}[breakable,colback=teal!4,colframe=teal!35!black,title={Before you move on}]
What does \mr{बादली} mean? (\textquotedblleft{}A bucket\textquotedblright{}.) Say it once more.
\end{tcolorbox}
\section[\texorpdfstring{\mr{दोरी} (dorī)}{दोरी (dorī)}]{\mr{दोरी} (dorī) — a rope}

\label{lesson:MR-C131-dori}



\begin{quote}
Before the new one: say the Marathi for a lock, then the Marathi for a bucket.
\end{quote}

\subsection*{You'll want to know: \mr{दोरी}}
\textbf{\mr{दोरी}} — \emph{dorī} — \textquotedblleft{}a rope\textquotedblright{}.

\begin{cousinweb}
\textbf{\mr{दोरी}} (\emph{dorī}) is a tadbhava, a word inherited from Sanskrit and reshaped by sound change. It comes from \emph{davara}, \textquotedblleft{}string\textquotedblright{} (CDIAL 6225).
\end{cousinweb}

\begin{grammarlens}[title={What you've built}]
One more word for this chapter.
\end{grammarlens}

\subsection*{Your turn}
\begin{itemize}
\raggedright
  \item \emph{Say it:} \emph{dorī}
  \item \emph{Say it:} \emph{dorī}, once more
  \item \emph{Say it:} say \emph{dorī}
  \item \emph{Recall:} say the Marathi for a lock, then the Marathi for a bucket, then say \emph{dorī} again
\end{itemize}

\begin{tcolorbox}[breakable,colback=teal!4,colframe=teal!35!black,title={Before you move on}]
What does \mr{दोरी} mean? (\textquotedblleft{}A rope\textquotedblright{}.) Say it once more.
\end{tcolorbox}
\section[\texorpdfstring{\mr{खिळा} (khiḷā)}{खिळा (khiḷā)}]{\mr{खिळा} (khiḷā) — a nail (to hammer)}

\label{lesson:MR-C131-khila}



\begin{quote}
Before the new one: say the Marathi for a bucket, then the Marathi for a rope.
\end{quote}

\subsection*{You'll want to know: \mr{खिळा}}
\textbf{\mr{खिळा}} — \emph{khiḷā} — \textquotedblleft{}a nail (to hammer)\textquotedblright{}.

\begin{cousinweb}
\textbf{\mr{खिळा}} (\emph{khiḷā}) belongs with Marathi \emph{khīḷ}, \textquotedblleft{}bolt, peg\textquotedblright{}, which CDIAL lists under Sanskrit \emph{khīla}, \textquotedblleft{}post\textquotedblright{}, a variant of \emph{kīla}, \textquotedblleft{}stake, peg\textquotedblright{} (CDIAL 3202). CDIAL thinks the variant shapes support a non-Aryan origin.
\end{cousinweb}

\begin{grammarlens}[title={What you've built}]
One more word for this chapter.
\end{grammarlens}

\subsection*{Your turn}
\begin{itemize}
\raggedright
  \item \emph{Say it:} \emph{khiḷā}
  \item \emph{Say it:} \emph{khiḷā}, once more
  \item \emph{Say it:} say \emph{khiḷā}
  \item \emph{Recall:} say the Marathi for a bucket, then the Marathi for a rope, then say \emph{khiḷā} again
\end{itemize}

\begin{tcolorbox}[breakable,colback=teal!4,colframe=teal!35!black,title={Before you move on}]
What does \mr{खिळा} mean? (\textquotedblleft{}A nail (to hammer)\textquotedblright{}.) Say it once more.
\end{tcolorbox}
